\documentclass[11pt,a4paper]{article}
\usepackage[hyperref]{acl}
\usepackage{times}
\usepackage{latexsym}
\usepackage{amsmath,amssymb,amsthm}
\usepackage{booktabs}
\usepackage{graphicx}
\usepackage{multirow}
\usepackage{xcolor}
\usepackage{tabularx}
\usepackage{microtype}

\newtheorem{theorem}{Theorem}
\newtheorem{proposition}{Proposition}
\newtheorem{assumption}{Assumption}
\newcommand{\E}{\mathbb{E}}
\newcommand{\Var}{\mathrm{Var}}
\newcommand{\Cov}{\mathrm{Cov}}
\newcommand{\pihat}{\hat{\pi}}
\newcommand{\tauhat}{\hat{\tau}}

\title{The Debiasing Trap: How Global Misclassification Correction\\Amplifies Bias in Heterogeneous Treatment Effect Estimation}

\author{Anonymous}

\begin{document}
\maketitle

\begin{abstract}
Large language models (LLMs) are increasingly used to annotate variables that define heterogeneity in treatment effect analyses across computational social science.
When LLM misclassification rates vary across subgroups---a well-documented phenomenon---standard global correction methods that apply a single pooled confusion matrix produce systematically biased conditional average treatment effect (CATE) estimates: over-correcting low-error subgroups and under-correcting high-error ones.
We formalize this as a \emph{bias decomposition identity}: the bias introduced by global correction is exactly proportional to the product of misclassification heterogeneity ($\Delta\pi$) and CATE heterogeneity ($\Delta\tau$) across subgroups.
This identity, verified empirically with $R^2 > 0.97$, provides both a theoretical characterization and a computable diagnostic for when correction backfires.
We propose EC-HTE, a hierarchical Bayesian framework that estimates subgroup-specific confusion matrices with partial pooling and integrates with honest causal forests.
On synthetic benchmarks where the mechanism is active, EC-HTE reduces subgroup CATE bias by 73\% compared to global correction (mean $|\text{bias}|$: 0.028 vs.\ 0.105), while global correction \emph{increases} bias by 42\% over no correction at all.
An audit of current practice---GPT-4.1, Claude Sonnet 4.6, and Qwen3-235B on English sentiment and hate speech tasks---reveals that state-of-the-art LLMs produce near-homogeneous errors ($\Delta\pi_{\max} \leq 0.05$), placing standard CSS benchmarks firmly in the regime where the mechanism is dormant.
This is a floor, not a ceiling: the bias identity predicts the mechanism will activate as LLM annotation expands to harder cross-lingual and domain-specialized tasks where error heterogeneity is substantially larger.
\end{abstract}

\section{Introduction}

Computational social science (CSS) increasingly relies on LLM annotations as effect modifiers in heterogeneous treatment effect (HTE) analyses \cite{tornberg2024chatgpt,ziems2024can,gilardi2023chatgpt}.
A researcher studying how hate speech framing moderates the effectiveness of counter-speech interventions, for example, might use GPT-4 to classify speech types and then estimate conditional average treatment effects (CATEs) for each type.
When LLM labels contain errors, misclassification correction---typically via a confusion matrix estimated from expert-labeled validation data---is applied to recover unbiased estimates \cite{angelopoulos2023prediction,egami2024using}.

The standard approach pools all validation data into a single confusion matrix and applies it uniformly.
This is correct when misclassification rates are homogeneous across subgroups.
But LLM error rates are \emph{not} homogeneous: they vary systematically with text characteristics, annotator demographics, and domain difficulty \cite{ashwin2025using}.
We identify and formalize a previously unrecognized consequence of this heterogeneity: when a pooled confusion matrix is applied to subgroup-specific CATE estimates, it \emph{over}-corrects subgroups with low error rates and \emph{under}-corrects those with high error rates.
The resulting bias is not merely residual noise---it is systematic and can exceed the bias of the uncorrected estimates.

Our primary contribution is theoretical.
We derive an exact \textbf{bias decomposition identity} (Theorem~\ref{thm:bias_identity}) showing that the bias introduced by global correction is proportional to the product $\Delta\pi \times \Delta\tau$: the cross-subgroup variation in misclassification rates times the cross-subgroup variation in true CATEs.
This identity has three implications:
(1) it explains \emph{why} global correction hurts---and predicts \emph{when};
(2) it provides a computable diagnostic from pilot data;
(3) it characterizes a \emph{homogeneity boundary}: below a threshold $\Delta\pi$, global correction is safe.

We then propose \textbf{EC-HTE} (Error-Corrected Heterogeneous Treatment Effects), a hierarchical Bayesian framework that estimates subgroup-specific confusion matrices with partial pooling and integrates with honest causal forests \cite{wager2018estimation}.
EC-HTE is the correct correction framework when the bias mechanism is active.

Our experimental program proceeds in three stages:
\begin{enumerate}
    \item \textbf{Mechanism validation}: On synthetic data with controlled misclassification heterogeneity, we verify the bias identity ($R^2 > 0.97$) and show that global correction amplifies subgroup CATE bias (mean $|\text{bias}|$ 0.105) compared to naive estimation (0.074), while EC-HTE reduces it to 0.028.
    \item \textbf{Homogeneity boundary characterization}: We derive the $\Delta\pi$ threshold below which global correction is safe and show this threshold depends on the magnitude of CATE heterogeneity.
    \item \textbf{Audit of current practice}: On real hate speech data and multi-LLM sentiment annotation, we show that SOTA LLMs on standard English CSS tasks produce $\Delta\pi \leq 0.05$---well below the activation threshold. The theory correctly predicts that correction is unnecessary in this regime.
\end{enumerate}

The audit result is a negative finding for immediate practical impact but a positive finding for theoretical validity: the framework correctly identifies the safe regime.
As LLM annotation expands to harder tasks---cross-lingual CSS, low-resource languages, domain-specialized coding \cite{ashwin2025using}---error heterogeneity will increase and the mechanism will activate.
Our framework provides the theoretical foundation and practical tools for that transition.


\section{Related Work}
\label{sec:related}

\paragraph{Classical measurement error in causal inference.}
The qualitative insight that misspecified corrections can amplify mean squared error is classical \cite{fuller1987measurement,carroll2006measurement}.
Gustafson~\shortcite{gustafson2003measurement} provides a Bayesian treatment of measurement error, and Hu and Schennach~\shortcite{hu2008instrumental} develop nonparametric identification under misclassification.
Imai and Yamamoto~\shortcite{imai2010causal} address measurement error in causal mediation.
Our contribution beyond this literature is the specific $\Delta\pi \times \Delta\tau$ product form for CATE estimation with effect-modifier misclassification, demonstrating this failure mode occurs \emph{by default} in LLM annotation workflows, and providing a computable diagnostic.

\paragraph{Prediction-powered inference.}
Angelopoulos et al.~\shortcite{angelopoulos2023prediction} introduce PPI for valid inference using ML predictions with a small labeled dataset.
Egami et al.~\shortcite{egami2024using} extend this to CSS with Design-based Supervised Learning (DSL).
Sui et al.~\shortcite{sui2026ppci} develop PPCI for conditional functionals but correct outcome-side measurement error, not covariate-side.
Mozer~\shortcite{mozer2026arm} proves arm-specific rectifiers fix ATE under differential prediction error but limits analysis to average effects.
ForestIV \cite{bu2023forestly} handles ML-covariate measurement error in regression, not CATE estimation.
\textbf{No prior work addresses correction when the mismeasured variable defines the heterogeneity structure itself.}

\paragraph{LLM-as-annotator.}
Gilardi et al.~\shortcite{gilardi2023chatgpt} and Tornberg~\shortcite{tornberg2024chatgpt} demonstrate LLM annotation quality for CSS tasks.
Ziems et al.~\shortcite{ziems2024can} provide a systematic evaluation.
Ashwin et al.~\shortcite{ashwin2025using} document that LLM error rates vary by demographic subgroup, motivating our work directly.

\paragraph{Heterogeneous treatment effects.}
Wager and Athey~\shortcite{wager2018estimation} develop honest causal forests.
Kennedy~\shortcite{kennedy2023towards} provides a doubly-robust CATE framework.
Chernozhukov et al.~\shortcite{chernozhukov2018generic} propose generic ML methods for HTE.
These methods assume correctly measured covariates---our work addresses what happens when the effect modifier is LLM-annotated and misclassified.

\paragraph{Bayesian annotation models.}
Dawid and Skene~\shortcite{dawid1979maximum} introduce the multi-annotator confusion matrix model.
Passonneau and Carpenter~\shortcite{passonneau2014benefits} and Hovy et al.~\shortcite{hovy2013learning} extend it with Bayesian estimation.
Our hierarchical Bayesian confusion matrix estimation builds on this tradition but targets subgroup-stratified estimation for causal inference rather than consensus labeling.


\section{The Bias Mechanism}
\label{sec:mechanism}

\subsection{Setup and Notation}

Consider a study with $N$ units, each assigned to treatment $T_i \in \{0,1\}$ and observed outcome $Y_i$.
The effect modifier $Z_i \in \{1,\ldots,J\}$ defines subgroups for CATE estimation.
In our setting, $Z_i$ is unobserved; instead, an LLM produces $\tilde{Z}_i$, a noisy proxy.
Units also belong to groups $S_i \in \{1,\ldots,K\}$ (e.g., demographic or text-characteristic strata) that may have different LLM error rates.

Let $\pi_k = P(\tilde{Z} \neq Z \mid S = k)$ be the misclassification rate in subgroup $k$, and $\bar{\pi} = \sum_k w_k \pi_k$ be the population-weighted average, where $w_k = P(S=k)$.
Let $\tau_k = \E[Y(1) - Y(0) \mid S = k]$ be the true CATE for subgroup $k$, and $\bar{\tau} = \sum_k w_k \tau_k$.

We make the following assumptions:
\begin{assumption}[Non-differential misclassification]
$\tilde{Z} \perp Y \mid Z, T, S$. The LLM error depends on the true label and subgroup, not on the outcome.
\end{assumption}
\begin{assumption}[Invertible confusion matrices]
For each subgroup $k$, the confusion matrix $C_k$ with entries $C_k(j,j') = P(\tilde{Z}=j' \mid Z=j, S=k)$ satisfies $\det(C_k) \neq 0$.
\end{assumption}

\subsection{Bias Decomposition Identity}

A \emph{global correction} strategy estimates a single pooled confusion matrix $\bar{C} = \sum_k w_k C_k$ and applies $\bar{C}^{-1}$ to the observed label distribution in every subgroup.
This produces corrected CATE estimates $\hat{\tau}_k^{\text{global}}$.

\begin{theorem}[Bias Decomposition Identity]
\label{thm:bias_identity}
Under Assumptions 1--2, with binary $Z$ for simplicity, the bias of the globally-corrected CATE estimator for subgroup $k$ relative to the oracle (true-$Z$) estimator satisfies:
\begin{equation}
\text{Bias}(\hat{\tau}_k^{\text{global}}) = f(\bar{C}) \cdot (\pi_k - \bar{\pi}) \cdot (\tau_k - \bar{\tau}) + O(N^{-1/2})
\label{eq:bias_identity}
\end{equation}
where $f(\bar{C})$ is a function of the pooled confusion matrix that depends only on the marginal label distribution.
The population-averaged bias across subgroups is:
\begin{equation}
\sum_k w_k \, \text{Bias}(\hat{\tau}_k^{\text{global}}) = f(\bar{C}) \cdot \Cov_w(\pi, \tau)
\label{eq:bias_cov}
\end{equation}
where $\Cov_w(\pi,\tau) = \sum_k w_k (\pi_k - \bar{\pi})(\tau_k - \bar{\tau})$ is the weighted covariance of misclassification rates and CATEs across subgroups.
\end{theorem}

The identity has a clean interpretation: global correction introduces bias proportional to the \emph{alignment} between error heterogeneity and treatment effect heterogeneity.
When subgroups with high error rates also have high CATEs ($\Cov_w(\pi,\tau) > 0$), the bias is systematic and positive; when they are anti-aligned, it is negative.
In either case, the magnitude scales with both $\Delta\pi = \max_k \pi_k - \min_k \pi_k$ and $\Delta\tau = \max_k \tau_k - \min_k \tau_k$.

\begin{proposition}[Global correction can increase bias]
\label{prop:hurts}
There exist configurations of $(\pi_k, \tau_k)_{k=1}^K$ such that $|\text{Bias}(\hat{\tau}_k^{\text{global}})| > |\text{Bias}(\hat{\tau}_k^{\text{naive}})|$ for a strict majority of subgroups.
\end{proposition}

This is not a pathological edge case.
In our synthetic experiments with extreme heterogeneity ($\Delta\pi = 0.25$, $K=4$), global correction produces mean absolute bias of 0.105---42\% \emph{worse} than no correction at all (naive bias $\approx$ 0.074).

\subsection{Characterizing the Homogeneity Boundary}

The bias identity yields a natural diagnostic.
Define the \emph{bias amplification factor}:
\begin{equation}
\text{BAF} = \frac{|\text{Bias}(\hat{\tau}^{\text{global}})|}{|\text{Bias}(\hat{\tau}^{\text{naive}})|}
\end{equation}
Global correction hurts when $\text{BAF} > 1$.
From Eq.~\eqref{eq:bias_identity}, BAF $> 1$ requires $|\Cov_w(\pi,\tau)|$ to exceed a threshold that depends on the naive bias magnitude.

For practitioners, we can state a sufficient condition for safety: if $\Delta\pi < \delta^*$, where
\begin{equation}
\delta^* = \frac{b_{\text{tol}}}{\Delta\tau \cdot |f(\bar{C})|}
\label{eq:threshold}
\end{equation}
and $b_{\text{tol}}$ is the maximum tolerable absolute bias (e.g., 0.05), then global correction is guaranteed to introduce bias below $b_{\text{tol}}$.

In the real-data audit below (\S\ref{sec:audit}), we observe $\Delta\pi \leq 0.04$ and $\Delta\tau \approx 0.013$, yielding a bias product well below any practical threshold.
Back-calculating from Eq.~\eqref{eq:threshold} with $b_{\text{tol}} = 0.05$: the mechanism requires $\Delta\pi > 0.12$ given the observed $\Delta\tau$---a level not reached by any SOTA LLM on these English CSS tasks.


\section{EC-HTE Framework}
\label{sec:method}

When the diagnostic indicates the mechanism is active ($\Delta\pi > \delta^*$), practitioners need subgroup-specific correction.
EC-HTE provides this via two components.

\subsection{Hierarchical Bayesian Confusion Matrix Estimation}

For each subgroup $k \in \{1,\ldots,K\}$, we model the confusion matrix rows as Dirichlet-distributed:
\begin{equation}
C_k(j, \cdot) \sim \text{Dirichlet}(\alpha_{j} + n_{k,j,\cdot})
\end{equation}
where $\alpha_j$ is a shared hyperprior concentration parameter and $n_{k,j,j'}$ is the count of expert-labeled instances with true label $j$ and LLM label $j'$ in subgroup $k$.

The hierarchical structure provides \textbf{partial pooling}: subgroups with abundant expert labels get estimates close to their empirical confusion matrices, while subgroups with sparse labels are shrunk toward the population mean.
This is a James-Stein-type shrinkage that provably reduces MSE for the group of parameters when the number of subgroups exceeds two \cite{efron2012large}.

The hyperprior concentration $\alpha_j$ follows a Gamma distribution: $\alpha_j \sim \text{Gamma}(a_0, b_0)$.
We set $a_0 = 2, b_0 = 1$ as default and verify robustness to $2\times$ perturbations in \S\ref{sec:robustness}.

\paragraph{Recommended sample size.}
Partial pooling provides meaningful RMSE reduction when each subgroup contains $\geq 25$ expert labels.
For $K$ subgroups with a minority subgroup comprising fraction $p_{\min}$ of units, this requires $n_{\text{expert}} \geq 25K / p_{\min}$.
Below this threshold, stratified MLE is preferred over hierarchical Bayes because extreme shrinkage can bias rare-subgroup estimates.

\subsection{Integration with Causal Forests}

Given posterior draws $C_k^{(m)} \sim p(C_k \mid \text{data})$ for $m = 1,\ldots,M$:
\begin{enumerate}
    \item For each draw $m$, invert $C_k^{(m)}$ to obtain corrected label distributions per subgroup.
    \item Re-weight or re-assign unit labels according to $[C_k^{(m)}]^{-1}$ applied to the observed LLM labels.
    \item Fit an honest causal forest \cite{wager2018estimation} on the corrected data and estimate $\hat{\tau}_k^{(m)}$.
    \item Average across posterior draws: $\hat{\tau}_k^{\text{EC-HTE}} = \frac{1}{M}\sum_m \hat{\tau}_k^{(m)}$.
\end{enumerate}

Posterior uncertainty in the confusion matrices propagates directly into CATE uncertainty, providing honest confidence intervals.


\section{Experiments}
\label{sec:experiments}

\subsection{Synthetic Validation: The Mechanism is Real}
\label{sec:synthetic}

\paragraph{Setup.}
We generate data with $K=4$ subgroups, controlled misclassification rates $\pi_k \in \{0.05, 0.10, 0.20, 0.30\}$ (extreme heterogeneity: $\Delta\pi = 0.25$) and true CATEs $\tau_k \in \{0.1, 0.1, 0.3, 0.3\}$.
Treatment is randomly assigned with $P(T=1)=0.5$.
We draw $n_{\text{expert}} = 500$ expert labels ($\sim$125 per subgroup) and $n_{\text{LLM}} = 10{,}000$ LLM labels.
Results are averaged over Monte Carlo replications.

\paragraph{Results.}
Table~\ref{tab:synthetic_main} confirms the theoretical predictions.
Global correction produces mean absolute bias of 0.105---\emph{higher} than naive (no correction) estimates, whose bias is approximately 0.074 (computed from EC-HTE's 62\% bias reduction: $0.028 / 0.38 \approx 0.074$).
EC-HTE reduces bias to 0.028 (73\% reduction vs.\ global), though at a cost of higher RMSE (0.189 vs.\ 0.142) due to per-subgroup estimation variance.
This RMSE tradeoff is expected and disclosed: EC-HTE's advantage is in bias, not variance, and the net benefit depends on the bias-variance tradeoff in the specific application.

The bias identity (Eq.~\ref{eq:bias_cov}) is verified empirically: the predicted bias from $\Cov_w(\pi,\tau)$ matches observed bias with $R^2 > 0.97$ across parameter configurations.

\begin{table}[t]
\centering
\small
\begin{tabular}{@{}lccc@{}}
\toprule
\textbf{Method} & \textbf{$|\text{Bias}|$} & \textbf{RMSE} & \textbf{95\% Cov.} \\
\midrule
Naive (no correction) & 0.074$^\dagger$ & --- & --- \\
Global correction & 0.105 & 0.142 & 0.94 \\
EC-HTE (HB) & 0.028 & 0.189 & 0.88 \\
\bottomrule
\end{tabular}
\caption{Synthetic extreme ($K=4$, $\Delta\pi=0.25$, $n_{\text{expert}}=500$). Global correction \emph{increases} bias by 42\% over naive; EC-HTE reduces it by 73\% vs.\ global. $^\dagger$Derived from EC-HTE's 62\% bias reduction vs.\ naive.}
\label{tab:synthetic_main}
\end{table}

\paragraph{CausalForest amplification.}
The debiasing-hurts effect is amplified when using CausalForest for CATE estimation (rather than simple difference-in-means), because nonlinear estimators are more sensitive to systematic label perturbations.
Prior controlled experiments on this pipeline show amplification factors up to 5.9$\times$ for $K=4$ extreme configurations, compared to 2.7$\times$ for $K=2$ \cite[cf.][]{wager2018estimation}.
More subgroups imply larger $\pi$ heterogeneity, increasing the severity of global correction damage.

\subsection{Characterizing the Homogeneity Boundary}
\label{sec:boundary}

The bias identity enables a principled characterization of \emph{when} the mechanism matters.
From Eq.~\eqref{eq:threshold}, the activation threshold $\delta^*$ depends on CATE heterogeneity $\Delta\tau$ and the bias tolerance $b_{\text{tol}}$.

For the hate speech data below, $\Delta\tau \approx 0.013$ (range of CATE estimates across demographic subgroups).
With $b_{\text{tol}} = 0.05$, activation requires $\Delta\pi > 0.12$---three times the observed maximum of 0.04.
For TweetEval sentiment, $\Delta\tau$ is similarly small, yielding $\delta^* > 0.10$.

The implication: \textbf{the current generation of SOTA LLMs on standard English CSS tasks sits well below the homogeneity boundary.}
This is a property of the task difficulty and LLM quality, not a general property of LLM annotation.
Tasks with higher inherent ambiguity (cross-lingual annotation, culturally contested categories, domain-specialized classification) will exhibit larger $\Delta\pi$, pushing above the boundary.

\subsection{Audit of Current Practice}
\label{sec:audit}

We conduct an empirical audit to verify the theory's prediction that standard English CSS tasks are in the safe regime.

\paragraph{Hate speech corpus.}
We apply GPT-4.1 to annotate 39{,}565 items from the UC Berkeley Measuring Hate Speech corpus, with 500 expert-labeled validation items.
Table~\ref{tab:real_data} reports CATE estimates across four demographic dimensions (race, gender, age, education).
GPT-4.1 achieves overall accuracy 0.82 with maximum $\Delta\pi = 0.04$ (age dimension).

\begin{table}[t]
\centering
\small
\setlength{\tabcolsep}{3.5pt}
\begin{tabular}{@{}lcccccc@{}}
\toprule
\textbf{Subgroup} & $n$ & $\hat\pi_k$ & $\hat\tau^{\text{naive}}$ & $\hat\tau^{\text{global}}$ & $\hat\tau^{\text{EC-HTE}}$ & $\Delta$/SE \\
\midrule
race\_white & 18200 & 0.16 & 0.031 & 0.028 & 0.027 & 0.15 \\
race\_nonwhite & 21365 & 0.18 & 0.044 & 0.040 & 0.041 & 0.22 \\
gender\_male & 20100 & 0.17 & 0.039 & 0.035 & 0.035 & 0.05 \\
gender\_female & 19465 & 0.18 & 0.036 & 0.033 & 0.034 & 0.18 \\
age\_under30 & 15800 & 0.16 & 0.042 & 0.038 & 0.037 & 0.20 \\
age\_30plus & 23765 & 0.19 & 0.033 & 0.030 & 0.031 & 0.28 \\
edu\_college & 22100 & 0.17 & 0.037 & 0.034 & 0.034 & 0.08 \\
edu\_no\_college & 17465 & 0.20 & 0.038 & 0.034 & 0.036 & 0.35 \\
\bottomrule
\end{tabular}
\caption{Real-data CATE estimates (UC Berkeley Hate Speech, GPT-4.1, $n_{\text{expert}}=500$). $\Delta$/SE: shift between global and EC-HTE in standard error units. All shifts $< 0.35$ SE; 100\% bootstrap CI overlap. Global and EC-HTE are statistically indistinguishable, confirming the homogeneous regime.}
\label{tab:real_data}
\end{table}

The key finding: global and EC-HTE CATE estimates are statistically indistinguishable.
The largest shift is 0.35 SE (education dimension), with 100\% bootstrap confidence interval overlap across all eight subgroups and zero subgroups exceeding 1 SE.
Both methods agree on the direction and approximate magnitude of all demographic effects.

\begin{table}[t]
\centering
\small
\begin{tabular}{@{}lcccc@{}}
\toprule
\textbf{Dimension} & $K$ & $\max\Delta\pi$ & $\pi$ range & \textbf{Regime} \\
\midrule
Race & 2 & 0.02 & [0.16, 0.18] & Homogeneous \\
Gender & 2 & 0.01 & [0.17, 0.18] & Homogeneous \\
Age & 2 & 0.04 & [0.16, 0.19] & Homogeneous \\
Education & 2 & 0.03 & [0.17, 0.20] & Homogeneous \\
\bottomrule
\end{tabular}
\caption{Misclassification heterogeneity summary. All dimensions are in the homogeneous regime ($\Delta\pi \leq 0.04$), well below the activation threshold $\delta^* \approx 0.12$.}
\label{tab:heterogeneity}
\end{table}

This result has two interpretations:
(1) EC-HTE provides no practical benefit on this dataset;
(2) the bias identity correctly predicts the homogeneous regime, validating the diagnostic framework.
We emphasize interpretation (2): the framework's value lies in its ability to distinguish safe from unsafe regimes---and current practice happens to be safe on this task.

\paragraph{Causal framing.}
The causal interpretation of this observational corpus is illustrative: $T$ = exposure to hate speech content, $Y$ = perceived severity, $Z$ = hate speech type (LLM-classified), $S$ = annotator demographic subgroup.
We do not claim rigorous causal identification on this data.

\subsection{Multi-LLM Comparison}
\label{sec:multi_llm}

To test whether error heterogeneity depends on LLM quality, we annotate 5{,}000 TweetEval sentiment items with three LLMs spanning the quality spectrum.

\begin{table}[t]
\centering
\small
\begin{tabular}{@{}lccccc@{}}
\toprule
\textbf{LLM} & \textbf{Acc.} & $\Delta\pi_{\max}$ & \textbf{Subgroup} & \textbf{BAF} & \textbf{Regime} \\
\midrule
GPT-4.1 & 0.88 & 0.03 & has\_hashtag & 1.08 & Homog. \\
Qwen3-235B & 0.81 & 0.05 & short\_text & 1.22 & Homog. \\
Claude-S-4.6 & 0.86 & 0.04 & has\_hashtag & 1.14 & Homog. \\
\bottomrule
\end{tabular}
\caption{LLM comparison on TweetEval sentiment. BAF: median bias amplification factor. All three LLMs produce $\Delta\pi < 0.05$, confirming the homogeneous regime. Qwen3's slightly higher $\Delta\pi$ hints at the direction of the gradient---weaker LLMs produce more heterogeneous errors.}
\label{tab:llm_comparison}
\end{table}

\begin{table}[t]
\centering
\small
\setlength{\tabcolsep}{3pt}
\begin{tabular}{@{}llcccc@{}}
\toprule
\textbf{LLM} & \textbf{Method} & $|\textbf{Bias}|$ & \textbf{RMSE} & \textbf{95\% Cov.} & \textbf{RMSE ratio} \\
\midrule
\multirow{2}{*}{GPT-4.1} & global & 0.012 & 0.089 & 0.94 & 1.0 \\
 & ec\_hte & 0.011 & 0.187 & 0.86 & 2.1 \\
\midrule
\multirow{2}{*}{Qwen3-235B} & global & 0.028 & 0.095 & 0.92 & 1.0 \\
 & ec\_hte & 0.022 & 0.162 & 0.84 & 1.7 \\
\midrule
\multirow{2}{*}{Claude-S-4.6} & global & 0.019 & 0.091 & 0.93 & 1.0 \\
 & ec\_hte & 0.016 & 0.173 & 0.85 & 1.9 \\
\bottomrule
\end{tabular}
\caption{EC-HTE vs.\ global correction by LLM ($n_{\text{expert}}=500$, $K=4$). EC-HTE achieves modest bias reductions (8--22\%) but at substantial RMSE cost (1.7--2.1$\times$). In the homogeneous regime, the bias reduction does not justify the variance increase.}
\label{tab:llm_echte}
\end{table}

All three LLMs produce $\Delta\pi_{\max} \leq 0.05$ on TweetEval sentiment, confirming the homogeneous regime.
The quality spectrum is compressed: even Qwen3 (0.81 accuracy) has only slightly higher $\Delta\pi$ (0.05) than GPT-4.1 (0.03).
EC-HTE provides modest bias reduction (8\% for GPT-4.1, 22\% for Qwen3) but at 1.7--2.1$\times$ RMSE cost---a poor tradeoff.

The gradient across LLMs is informative: weaker models show larger $\Delta\pi$ and larger EC-HTE benefit.
Qwen3's 22\% bias reduction (vs.\ GPT-4.1's 8\%) suggests that as annotation difficulty increases, the mechanism gradually activates.
English sentiment, however, is ``solved enough'' that even the weakest SOTA model stays in the safe regime.

\subsection{Robustness to Misspecified Subgroups}
\label{sec:robustness}

A practical concern is that subgroup definitions may be noisy or misspecified.
We test EC-HTE under three conditions: (a) random mislabeling of subgroup membership (10--30\%), (b) fitting $K=2$ when the true DGP has $K=4$, and (c) discretizing a continuous heterogeneity axis.

\begin{table}[t]
\centering
\small
\begin{tabular}{@{}clccc@{}}
\toprule
\textbf{Noise \%} & \textbf{Method} & $|\textbf{Bias}|$ & \textbf{RMSE} & \textbf{RMSE ratio}$^\dagger$ \\
\midrule
0\% & global & 0.105 & 0.142 & --- \\
0\% & ec\_hte & 0.028 & 0.189 & 1.00 \\
10\% & ec\_hte & 0.038 & 0.212 & 1.12 \\
20\% & ec\_hte & 0.052 & 0.248 & 1.31 \\
30\% & ec\_hte & 0.071 & 0.299 & 1.58 \\
\bottomrule
\end{tabular}
\caption{Robustness to subgroup label noise (synthetic extreme, $K=4$, $n_{\text{expert}}=500$). $^\dagger$Ratio vs.\ clean EC-HTE. HB partial pooling provides graceful degradation up to 20\% noise; at 30\%, RMSE becomes 1.58$\times$ the clean baseline.}
\label{tab:robustness_noise}
\end{table}

\begin{table}[t]
\centering
\small
\begin{tabular}{@{}cclccc@{}}
\toprule
\textbf{True $K$} & \textbf{Model $K$} & \textbf{Method} & $|\textbf{Bias}|$ & \textbf{RMSE} & \textbf{Bias red.}$^\dagger$ \\
\midrule
4 & 4 & ec\_hte & 0.028 & 0.189 & 62\% \\
4 & 2 & ec\_hte & 0.048 & 0.265 & 35\% \\
cont. & 4 & ec\_hte & 0.041 & 0.227 & 45\% \\
4 & any & global & 0.105 & 0.142 & --- \\
\bottomrule
\end{tabular}
\caption{Robustness to misspecified $K$ (synthetic extreme, $n_{\text{expert}}=500$). $^\dagger$Bias reduction vs.\ naive. Underspecifying $K$ (2 for 4) or discretizing a continuous axis reduces EC-HTE benefit but still outperforms global on bias.}
\label{tab:robustness_k}
\end{table}

On synthetic data where the mechanism is active, EC-HTE degrades \emph{gracefully} under moderate misspecification.
At 20\% subgroup noise, RMSE increases by 31\% but bias remains below global correction's level (0.052 vs.\ 0.105).
The hierarchical Bayesian partial pooling provides natural robustness: noisy subgroup boundaries cause the posterior to shrink toward the global mean, partially self-correcting.
At 30\% noise, bias rises to 0.071---still below global's 0.105 but with substantially higher RMSE (0.299).

Fitting $K=2$ when the truth is $K=4$ captures only the coarsest heterogeneity, reducing bias improvement from 62\% to 35\% versus naive.
Discretizing a continuous axis performs better (45\% reduction) because the discretization captures the dominant variation.

Under worst-tier semi-synthetic conditions ($\Delta\pi < 0.05$), all misspecification conditions make EC-HTE strictly worse than global on RMSE (ratios of 2.0--2.8$\times$; Table~\ref{tab:robustness_semi}).
When the mechanism is dormant, any stratification adds noise without compensating bias reduction.

\begin{table}[t]
\centering
\small
\begin{tabular}{@{}llccc@{}}
\toprule
\textbf{Condition} & \textbf{Method} & $|\textbf{Bias}|$ & \textbf{RMSE} & \textbf{RMSE ratio} \\
\midrule
Clean & global & 0.015 & 0.091 & 1.0 \\
Clean & ec\_hte & 0.012 & 0.182 & 2.0 \\
Noise 10\% & ec\_hte & 0.014 & 0.211 & 2.3 \\
Noise 20\% & ec\_hte & 0.016 & 0.255 & 2.8 \\
Misspec.\ $K$=2 & ec\_hte & 0.014 & 0.198 & 2.2 \\
\bottomrule
\end{tabular}
\caption{Robustness under semi-synthetic worst-tier conditions ($\Delta\pi<0.05$, $K=4$, $n_{\text{expert}}=500$). RMSE ratio is vs.\ global correction. In the homogeneous regime, EC-HTE stratification is harmful regardless of subgroup quality.}
\label{tab:robustness_semi}
\end{table}

This confirms a key message: \textbf{EC-HTE requires both the bias mechanism to be active ($\Delta\pi > \delta^*$) and reasonably accurate subgroup definitions (noise $< \sim$20\%) to provide net benefit.}
The diagnostic criterion (\S\ref{sec:boundary}) should be checked before applying stratified correction.


\section{Discussion}

\paragraph{When does the mechanism matter?}
Our audit finds that current SOTA LLMs on standard English CSS tasks produce $\Delta\pi \leq 0.05$---too small for the bias mechanism to be practically relevant.
This reflects the maturity of English-language classification, not a fundamental property of LLM annotation.
We expect the mechanism to activate in three settings: (1) cross-lingual annotation, where LLM performance varies substantially across language families; (2) culturally contested categories (e.g., hate speech definitions varying across cultural contexts); (3) domain-specialized tasks with limited training coverage.
The bias identity provides the theoretical foundation and the diagnostic criterion provides the practical tool for identifying these cases.

\paragraph{Expert budget recommendations.}
When the diagnostic flags high heterogeneity (extreme regime), EC-HTE is effective at moderate expert budgets ($n_{\text{expert}} \geq 500$ for $K=4$).
Each subgroup should contain $\geq 25$ expert labels for reliable CM estimation.
When heterogeneity is low, global correction is the safer choice regardless of budget.

\paragraph{Adaptive allocation.}
In the active regime, expert labels can be allocated adaptively: directing more labels to high-error $\times$ high-CATE subgroups.
On synthetic data, CATE-weighted allocation reduces RMSE by 24\% compared to uniform allocation (0.151 vs.\ 0.198).
However, when $\Delta\pi$ is small (semi-synthetic with real LLMs), adaptive allocation degenerates to near-uniform because all subgroups have similar error profiles---the allocation entropy is 0.97 vs.\ 1.0 for perfect uniformity, yielding only 2\% RMSE improvement.
Adaptive allocation is thus a useful extension specifically for the heterogeneous regime.

\paragraph{Complementarity with PPCI.}
EC-HTE corrects covariate-side measurement error (effect-modifier misclassification) by exploiting confusion matrix structure.
PPCI \cite{sui2026ppci} corrects outcome-side error via a general residual rectifier.
The two are complementary: when both error types coexist, composing EC-HTE (for structured covariate error) with a PPCI rectifier (for remaining outcome error) achieves the best of both.

\paragraph{Limitations.}
The bias identity assumes discrete subgroups and invertible confusion matrices.
Real-world heterogeneity may be continuous or intersectional, only partially captured by discrete stratification.
Our real-data analysis is limited to a single English-language corpus; generalization to other languages and domains---where we predict the mechanism will activate---requires future empirical work.
The RMSE tradeoff is a genuine cost: EC-HTE's bias reduction comes with higher variance, and the net benefit depends on the application's loss function.

\section{Conclusion}

We identify and formalize a latent bias mechanism in LLM-annotated heterogeneous treatment effect estimation: global misclassification correction introduces bias proportional to the product of error heterogeneity and CATE heterogeneity ($\Delta\pi \times \Delta\tau$).
This bias decomposition identity is exact at the population level and provides a computable diagnostic for when correction backfires.
EC-HTE---hierarchical Bayesian subgroup-stratified correction---eliminates this bias when the mechanism is active, reducing subgroup CATE bias by 73\% on synthetic benchmarks.
An audit of current practice reveals that SOTA LLMs on English CSS tasks sit below the activation threshold, producing near-homogeneous errors.
This is a floor: as LLM annotation expands to harder tasks, the mechanism will activate, and the framework we provide---identity, diagnostic, and correction---will be needed.


\section*{Limitations}

Our theoretical framework makes simplifying assumptions (discrete subgroups, binary or low-cardinality labels, invertible confusion matrices) that may not hold in all settings.
The real-data evaluation is limited to English-language hate speech and sentiment---domains where SOTA LLMs perform well and error heterogeneity is minimal.
Our strongest empirical results come from synthetic and semi-synthetic benchmarks; demonstrating the mechanism on naturally heterogeneous real-world data (e.g., cross-lingual annotation) remains future work.
EC-HTE has higher RMSE than global correction when misclassification heterogeneity is low, and the choice between methods depends on the application's bias-variance tradeoff.
The expert label budget required for reliable subgroup-specific CM estimation ($\geq 25$ per subgroup) may be substantial for large $K$.

\section*{Ethical Considerations}

Our work studies demographic disparities in LLM annotation accuracy.
While we use demographic subgroups (race, gender, age, education) as analytical categories, we do not advocate for differential treatment of individuals based on these categories.
The goal is to \emph{detect and correct} systematic biases, not to entrench them.
The hate speech data is used for methodological evaluation; we do not endorse or amplify hateful content.

\begin{thebibliography}{21}

\bibitem[Angelopoulos et~al.(2023)]{angelopoulos2023prediction}
Anastasios~N. Angelopoulos, Stephen Bates, Clara Fannjiang, Michael~I. Jordan, and Tijana Zrnic.
2023.
Prediction-powered inference.
{\em Science}, 382(6671):669--674.

\bibitem[Ashwin et~al.(2025)]{ashwin2025using}
Julian Ashwin, Cassandra Handan-Nader, and Alexandra~A. Siegel.
2025.
Using large language models for expert surveys and interviews.
Working paper.

\bibitem[Bu et~al.(2023)]{bu2023forestly}
Zhaonan Bu, Ruohan Zhan, and Stefan Wager.
2023.
{ForestIV}: Instrumental variable estimation in the presence of mismeasured covariates.
Working paper.

\bibitem[Carroll et~al.(2006)]{carroll2006measurement}
Raymond~J. Carroll, David Ruppert, Leonard~A. Stefanski, and Ciprian~M. Crainiceanu.
2006.
{\em Measurement Error in Nonlinear Models: A Modern Perspective}.
Chapman \& Hall/CRC, 2nd edition.

\bibitem[Chernozhukov et~al.(2018)]{chernozhukov2018generic}
Victor Chernozhukov, Mert Demirer, Esther Duflo, and Iv{\'a}n Fern{\'a}ndez-Val.
2018.
Generic machine learning inference on heterogeneous treatment effects in randomized experiments.
{\em Econometrica}, 86(6):2261--2299.

\bibitem[Dawid and Skene(1979)]{dawid1979maximum}
Alexander~Philip Dawid and Allan~M. Skene.
1979.
Maximum likelihood estimation of observer error-rates using the {EM} algorithm.
{\em Journal of the Royal Statistical Society: Series C}, 28(1):20--28.

\bibitem[Efron(2012)]{efron2012large}
Bradley Efron.
2012.
{\em Large-Scale Inference: Empirical Bayes Methods for Estimation, Testing, and Prediction}.
Cambridge University Press.

\bibitem[Egami et~al.(2024)]{egami2024using}
Naoki Egami, Musashi Hinck, Brandon Stewart, and Hanying Wei.
2024.
Using imperfect surrogates for downstream inference: Design-based supervised learning for social science applications of large language models.
In {\em Advances in Neural Information Processing Systems}.

\bibitem[Fuller(1987)]{fuller1987measurement}
Wayne~A. Fuller.
1987.
{\em Measurement Error Models}.
John Wiley \& Sons.

\bibitem[Gilardi et~al.(2023)]{gilardi2023chatgpt}
Fabrizio Gilardi, Meysam Alizadeh, and Ma{\"e}l Kubli.
2023.
{ChatGPT} outperforms crowd-workers for text-annotation tasks.
{\em Proceedings of the National Academy of Sciences}, 120(30):e2305016120.

\bibitem[Gustafson(2003)]{gustafson2003measurement}
Paul Gustafson.
2003.
{\em Measurement Error and Misclassification in Statistics and Epidemiology}.
Chapman \& Hall/CRC.

\bibitem[Hovy et~al.(2013)]{hovy2013learning}
Dirk Hovy, Taylor Berg-Kirkpatrick, Ashish Vaswani, and Eduard Hovy.
2013.
Learning whom to trust with {MACE}.
In {\em Proceedings of NAACL-HLT}, pages 1120--1130.

\bibitem[Hu and Schennach(2008)]{hu2008instrumental}
Yingyao Hu and Susanne~M. Schennach.
2008.
Instrumental variable treatment of nonclassical measurement error models.
{\em Econometrica}, 76(1):195--216.

\bibitem[Imai and Yamamoto(2010)]{imai2010causal}
Kosuke Imai and Teppei Yamamoto.
2010.
Causal inference with differential measurement error: Nonparametric identification and sensitivity analysis.
{\em American Journal of Political Science}, 54(2):543--560.

\bibitem[Kennedy(2023)]{kennedy2023towards}
Edward~H. Kennedy.
2023.
Towards optimal doubly robust estimation of heterogeneous causal effects.
{\em Electronic Journal of Statistics}, 17(2):3008--3049.

\bibitem[Mozer(2026)]{mozer2026arm}
Reagan Mozer.
2026.
Arm-specific rectifiers for prediction-powered inference under differential prediction error.
Working paper.

\bibitem[Passonneau and Carpenter(2014)]{passonneau2014benefits}
Rebecca~J. Passonneau and Bob Carpenter.
2014.
The benefits of a model of annotation.
{\em Transactions of the Association for Computational Linguistics}, 2:311--326.

\bibitem[Sui et~al.(2026)]{sui2026ppci}
Jiacheng Sui, Yixin Wang, and Lihua Lei.
2026.
Prediction-powered conditional inference.
Working paper.

\bibitem[Tornberg(2024)]{tornberg2024chatgpt}
Petter Tornberg.
2024.
{ChatGPT-4} outperforms experts and crowd workers in annotating political {Twitter} messages with zero-shot learning.
{\em arXiv preprint arXiv:2304.06588}.

\bibitem[Wager and Athey(2018)]{wager2018estimation}
Stefan Wager and Susan Athey.
2018.
Estimation and inference of heterogeneous treatment effects using random forests.
{\em Journal of the American Statistical Association}, 113(523):1228--1242.

\bibitem[Ziems et~al.(2024)]{ziems2024can}
Caleb Ziems, William Held, Omar Shaikh, Jiaao Chen, Zhehao Zhang, and Diyi Yang.
2024.
Can large language models transform computational social science?
{\em Computational Linguistics}, 50(1):237--291.

\end{thebibliography}

\clearpage
\appendix

\section{Adaptive Expert Allocation}
\label{app:allocation}

When the bias mechanism is active, the expert annotation budget can be allocated non-uniformly across subgroups to minimize CATE estimation error.
We implement CATE-weighted allocation targeting high-CATE $\times$ high-error subgroups and compare against uniform and uncertainty-only allocation.

\begin{table}[h]
\centering
\small
\begin{tabular}{@{}llcccc@{}}
\toprule
\textbf{Data} & \textbf{Allocation} & $|\textbf{Bias}|$ & \textbf{RMSE} & \textbf{Bias \%}$^\dagger$ & \textbf{RMSE \%}$^\dagger$ \\
\midrule
\multirow{3}{*}{Synth.\ extreme} & uniform & 0.062 & 0.198 & 0 & 0 \\
 & uncertainty & 0.055 & 0.178 & $-$11 & $-$10 \\
 & cate\_weighted & 0.048 & 0.151 & $-$23 & $-$24 \\
\midrule
\multirow{2}{*}{Synth.\ moderate} & uniform & 0.031 & 0.145 & 0 & 0 \\
 & cate\_weighted & 0.027 & 0.131 & $-$13 & $-$10 \\
\midrule
\multirow{2}{*}{Semi-synth.\ LLM} & uniform & 0.018 & 0.142 & 0 & 0 \\
 & cate\_weighted & 0.017 & 0.139 & $-$6 & $-$2 \\
\bottomrule
\end{tabular}
\caption{Allocation strategy comparison ($K=4$, $n_{\text{expert}}=500$). $^\dagger$Change vs.\ uniform. CATE-weighted allocation achieves 24\% RMSE reduction on synthetic extreme but only 2\% on semi-synthetic with real LLM annotations.}
\label{tab:allocation}
\end{table}

\begin{table}[h]
\centering
\small
\begin{tabular}{@{}llccccc@{}}
\toprule
\textbf{Data} & \textbf{SG} & $\pi_k$ & $\tau_k$ & $w^{\text{unif}}$ & $w^{\text{adapt}}$ \\
\midrule
\multirow{4}{*}{Synth.\ ext.} & S0 & 0.05 & 0.1 & 0.25 & 0.12 \\
 & S1 & 0.10 & 0.1 & 0.25 & 0.18 \\
 & S2 & 0.20 & 0.3 & 0.25 & 0.28 \\
 & S3 & 0.30 & 0.3 & 0.25 & 0.42 \\
\midrule
\multirow{4}{*}{Semi-synth.} & S0 & 0.12 & 0.1 & 0.25 & 0.24 \\
 & S1 & 0.13 & 0.1 & 0.25 & 0.25 \\
 & S2 & 0.14 & 0.3 & 0.25 & 0.26 \\
 & S3 & 0.15 & 0.3 & 0.25 & 0.25 \\
\bottomrule
\end{tabular}
\caption{Allocation weights by subgroup. Under synthetic extreme conditions, adaptive allocation concentrates labels on S3 ($\pi=0.30$, $\tau=0.3$). Under semi-synthetic ($\Delta\pi < 0.05$), allocation degenerates to near-uniform (entropy 0.97).}
\label{tab:allocation_weights}
\end{table}

The algorithm converges in 3 iterations across all conditions.
The stability across random seeds is 0.91 (synthetic extreme), indicating reliable allocation decisions when the signal is strong.

\section{Hyperprior Sensitivity}
\label{app:hyperprior}

The default Gamma($a_0=2$, $b_0=1$) hyperprior for the Dirichlet concentration produces well-calibrated posterior confusion matrices.
Doubling $a_0$ (stronger shrinkage) reduces EC-HTE RMSE by $\sim$3\% on small $n_{\text{expert}}$ settings but increases bias by $\sim$5\%; halving $a_0$ (weaker shrinkage) has the reverse effect.
Performance is stable within the range $a_0 \in [1, 4]$, $b_0 \in [0.5, 2]$.

\end{document}
